% !TeX encoding = UTF-8
% !TeX program = xelatex
% !TeX root = GMCM-main.tex
% 主模板：examples/example-color.tex；功能与章节组织参考 CUMCM 模板。
% natbib 已由 gmcmthesis 加载，须在文档类之前传递新增选项。
\PassOptionsToPackage{sort&compress}{natbib}
\documentclass[bwprint,fontset=windows]{gmcmthesis}

% ==================== 保留 Paper 的版式 ====================
% 字体、字号、页边距、行距、标题、图表标题、封面和页码由 gmcmthesis.cls 管理。
% ctex、amsmath、amssymb、bm、graphicx、array、booktabs、tabularx、longtable、
% multirow、enumitem、caption、xcolor、listings、natbib、hyperref 已由类加载。
\hypersetup{hidelinks,pdfstartview=Fit,breaklinks=true}

% ==================== example-color 原有功能 ====================
\usepackage[framemethod=TikZ]{mdframed}
\usepackage{subfig} % 保留 \subfloat；不要同时加载 subcaption。
\usepackage{siunitx}
\usepackage{colortbl}
\definecolor{color1}{rgb}{0.78,0.88,0.99}
\definecolor{color2}{rgb}{0.36,0.62,0.84}
\definecolor{color3}{rgb}{0.88,0.92,0.96}
\definecolor{color4}{rgb}{0.96,0.97,0.98}
\newcommand{\red}[1]{\textcolor{red}{#1}}
\newcommand{\blue}[1]{\textcolor{blue}{#1}}

\usepackage[noend]{algpseudocode}
\usepackage{algorithm}
\floatname{algorithm}{算法}
\renewcommand{\algorithmicrequire}{\textbf{输入:}}
\renewcommand{\algorithmicensure}{\textbf{输出:}}

% ==================== 从 CUMCM 合并的内容功能 ====================
\usepackage{float}
\usepackage[figuresright]{rotating}
\usepackage{pdflscape}
\usepackage{afterpage}
\usepackage{multicol}
% 保留 Paper 的图片路径，尤其封面 logo.pdf、title.pdf 的 figures/。
\graphicspath{{figures/}{texfile/figures/}{figure/}{pictures/}%
  {picture/}{pic/}{pics/}{image/}{images/}}

% 具名代码样式：仅对显式 style=Python / style=MATLAB 的代码生效，
% 不覆盖 Paper 原有的 Matlab、Python 环境。
\lstdefinestyle{GMCMCode}{
  numbers=left,
  numberstyle=\zihao{-5}\ttfamily,
  numbersep=8pt,
  showspaces=false,
  showstringspaces=false,
  showtabs=false,
  captionpos=t,
  frame=lrtb,
  breaklines=true,
  columns=fixed,
  basewidth=0.5em,
  tabsize=4,
  xleftmargin=2em,
  basicstyle=\zihao{-5}\ttfamily\color{black}
}
\lstdefinestyle{Python}{
  style=GMCMCode,
  language=Python,
  keywordstyle=\color{blue}\bfseries\itshape,
  stringstyle=\color{magenta},
  commentstyle=\color{red}
}
\lstdefinestyle{MATLAB}{
  style=GMCMCode,
  language=Matlab,
  keywordstyle=\color{green!40!black}\bfseries\itshape,
  stringstyle=\color{purple},
  commentstyle=\color{gray}
}

% ==================== 论文信息：按实际参赛信息填写 ====================
\title{全国研究生数学建模竞赛论文标题}
\baominghao{\hspace{6em} 20260900001}
\schoolname{\hspace{5.8em} 学校名称填写}
\membera{\hspace{6em} 成员A}
\memberb{\hspace{6em} 成员B}
\memberc{\hspace{6em} 成员C}

\begin{document}
\maketitle
% !TeX root = ../GMCM-main.tex
\begin{abstract}
本文档是待填写的论文模板。此处填写研究背景、需要解决的问题、采用的主要方法以及核心结果；正文中的公式、图表和程序均为排版示例，应在写作时替换为实际内容。

\textbf{针对问题一}，此处概括模型的建立思路、求解方法及主要结果。

\textbf{针对问题二}，此处概括模型的改进、约束条件及主要结果。

\textbf{针对问题三}，此处概括模型的应用、验证方法及主要结果。

最后，此处概括误差或敏感性分析、模型的适用范围与主要结论。

\keywords{关键词一\quad 关键词二\quad 关键词三\quad 关键词四}
\end{abstract}

\pagestyle{plain}

% 默认从摘要进入正文；需要目录时再启用下面两行。
% \maketoc
% \clearpage

% !TeX root = ../GMCM-main.tex
\section{问题重述}
\subsection{问题背景}
在具身智能、自然人机交互和智能服务系统等领域，情感分析构成了不可或缺的关键支撑。在真实人机交互过程中，人的情绪并非由单一的形态信号进行表达，而是同时通过语言内容、语言韵律、面部神态以及身体动作等多个信号进行协同体现的。因此，仅依赖文本、语音或视觉中的某一模态进行情感判断，往往难以准确刻画复杂场景下的真实情感状态。多模态信息联合建模能够充分利用不同模态之间的互补性与一致性，是突破单模态的识别局限，进而提升复杂场景下对情绪理解的准确性和稳定性。
然而，在实际情景下，采集过程中更容易受到环境噪声、画面遮挡、设备故障和语音转写失败等因素影响，致使某一模态在局部时间段内缺失。若采用常规固定权重进行模态融合，可能会把缺失值误认为有效信息，进而造成预测结果性能的下降。与此同时，若仅给出情感类别或情感强度而不说明判断依据，则会降低模型在智能交互场景中的可复核性。基于上述背景，本题以CMU-MOSEI英文多模态情感数据为基础，针对进行多模态特征的提取与时序对齐、缺失性预测和可解释性情感这三方面进行研究。
\subsection{问题提出}
结合题目要求，需要解决以下问题：
\begin{enumerate}
  \item 问题一：填写需要求解的对象、条件与目标。
  \item 问题二：填写需要分析或优化的内容。
  \item 问题三：填写需要验证或推广的内容。
\end{enumerate}

% !TeX root = ../GMCM-main.tex
\section{问题分析}
\subsection{问题一分析}
此处说明问题一的特点、关键变量、模型选择依据和求解思路。

\subsection{问题二分析}
此处说明问题二与前一问题的联系、目标函数及主要约束。

\subsection{问题三分析}
此处说明问题三的验证方法、适用条件和需要比较的方案。

% !TeX root = ../GMCM-main.tex
\section{模型的假设}
根据实际问题补充合理假设，并说明假设的依据及影响：
\begin{enumerate}
  \item 假设一：待填写。
  \item 假设二：待填写。
  \item 假设三：待填写。
\end{enumerate}

% 符号表与标题一起排版，避免浮动到下一节之后。
\Needspace{10\baselineskip}
\section{符号说明}
% 标准三线表：仅顶线、表头分隔线、底线，不加竖线和逐行边框。
% 浅色只用于表头背景；列格式和行高只在本表内生效。
\begin{table}[H]
  \centering
  \renewcommand{\arraystretch}{1.2}
  \begin{tabularx}{0.9\textwidth}{>{\centering\arraybackslash}p{0.15\textwidth}X}
    \toprule
    \rowcolor{color3!45}\textbf{符号} & \textbf{意义} \\
    \midrule
    $x_i$ & 第 $i$ 个输入变量（示例） \\
    $y_i$ & 第 $i$ 个观测值（示例） \\
    $\hat{y}_i$ & 第 $i$ 个预测值（示例） \\
    $n$ & 样本数量（示例） \\
    \bottomrule
  \end{tabularx}
\end{table}

% !TeX root = ../GMCM-main.tex
\section{模型的建立与求解}
本节保留公式、图表、算法和引用的用法示例，具体模型与数据应替换为实际研究内容。
各问题统一按“模型建立、求解设置、结果、验证与解释”组织；以下填写提示不代表已经完成相应计算或检验。

\subsection{数据的预处理}
\begin{enumerate}
  \item 描述数据来源、变量含义与单位，明确缺失值和异常值的处理依据。
  \item 说明归一化、特征提取或样本划分方法，给出处理后的数据规模。
\end{enumerate}
单位排版示例：长度为 \SI{120}{\centi\metre}，时间为 \SI{10}{\second}。

\subsection{问题一的模型建立与求解}
\subsubsection{模型建立}
此处定义问题一的变量、目标与约束，说明模型假设及其与实际问题的对应关系。

以线性拟合为排版示例，令预测值为
\begin{equation}
  \hat{y}_i = a x_i + b,\qquad i=1,2,\ldots,n.
  \label{eq:prediction}
\end{equation}
相应的最小二乘目标可写为
\begin{equation}
  \min_{a,b}\sum_{i=1}^{n}\bigl(y_i-a x_i-b\bigr)^2.
  \label{eq:least-squares}
\end{equation}
式~\eqref{eq:prediction} 和式~\eqref{eq:least-squares} 展示公式编号及交叉引用。

\subsubsection{求解设置}
此处说明计算所用的数据范围、算法、软件版本、参数取值及其依据。若涉及迭代或随机计算，还需记录初始条件、停止准则、数值容差或随机种子，以便复现。

算法~\ref{alg:fit} 展示伪代码的输入、输出和计算步骤。
\begin{algorithm}[htbp]
  \caption{线性拟合示例}
  \label{alg:fit}
  \begin{algorithmic}[1]
    \Require 数据点 $\{(x_i,y_i)\}_{i=1}^{n}$，且 $x_i$ 不全相同
    \Ensure 拟合系数 $a,b$ 及预测值 $\hat{y}_i$
    \State 构造设计矩阵 $X=[\bm{x},\bm{1}]$ 和观测向量 $\bm{y}$
    \State 使用 QR 分解求解最小二乘问题 $X\bm{\beta}\approx\bm{y}$
    \State 从 $\bm{\beta}$ 读取 $a,b$
    \For{$i=1,2,\ldots,n$}
      \State $\hat{y}_i\gets a x_i+b$
    \EndFor
    \State \Return $a,b,\{\hat{y}_i\}_{i=1}^{n}$
  \end{algorithmic}
\end{algorithm}

\subsubsection{结果}
此处给出问题一的核心计算结果、单位和对应条件，结合图表直接回答题目，并区分计算结果与据此得到的结论。

表~\ref{tab:sample} 为三线表示例，数值仅用于演示排版。
\begin{table}[htbp]
  \centering
  \caption{示例数据（非实际计算结果）}
  \label{tab:sample}
  \renewcommand{\arraystretch}{1.2}
  % 与符号表统一：三条黑色横线、无竖线，仅表头使用浅色背景。
  \begin{tabular}{ccc}
    \toprule
    \rowcolor{color3!45}\textbf{样本编号} & \textbf{输入} $x_i$ & \textbf{观测值} $y_i$ \\
    \midrule
    1 & 1 & 3 \\
    2 & 2 & 5 \\
    3 & 3 & 7 \\
    \bottomrule
  \end{tabular}
\end{table}

图~\ref{fig:demo} 沿用 Paper 的示例图片及子图写法。
\begin{figure}[htbp]
  \centering
  \subfloat[示例图一]{\includegraphics[width=0.40\textwidth]{image1}}\qquad
  \subfloat[示例图二]{\includegraphics[width=0.40\textwidth]{image2}}
  \caption{并列子图示例（非实际研究结果）}
  \label{fig:demo}
\end{figure}

\subsubsection{验证与解释}
此处根据问题性质选择验证方法，例如独立数据检验、基准方法比较、约束可行性检查或残差分析。说明检验依据、评价指标和实际检验结果，解释这些证据如何支持本问结论，同时交代结果的适用条件与局限。

\subsection{问题二的模型建立与求解}
\subsubsection{模型建立}
此处定义问题二的变量、目标函数和约束条件。若沿用问题一的模型，应说明保留的部分、需要调整的假设以及新增条件。

\subsubsection{求解设置}
此处说明问题二采用的数据、求解算法、参数及计算精度；复用前文设置时明确引用，只展开差异。根据实际方法补充初值、停止准则、随机种子或计算资源。

\subsubsection{结果}
此处给出问题二的核心结果及对应条件，并直接回答题目。若比较多个方案，统一比较条件、指标和单位，说明差异的大小及其实际含义。

\subsubsection{验证与解释}
此处选择与问题二目标相符的检验，例如与基准方案比较、检查约束是否满足，或在不同数据与条件下复核结果。报告检验结果，解释方案有效的原因、可能失效的情形及结论适用范围。

\subsection{问题三的模型建立与求解}
\subsubsection{模型建立}
此处说明问题三的研究目标、变量关系和约束条件，明确与前两问的联系；若引入新模型或推广原模型，应给出相应依据和适用假设。

\subsubsection{求解设置}
此处说明数据或场景的选取依据、算法与参数设置，以及必要的实验或计算步骤。需要重复实验时，说明重复次数和结果汇总方法。

\subsubsection{结果}
此处按题目要求汇总问题三的主要结果，给出量化指标、单位和适用条件。若结果随场景变化，应说明各场景的结论及其差异。

\subsubsection{验证与解释}
此处说明如何检验问题三的结果，并报告实际验证证据。结合模型机制解释主要现象，讨论结果对数据、参数或场景变化的依赖，以及结论能够推广到何种范围。

% !TeX root = ../GMCM-main.tex
\section{误差分析}
\subsection{误差来源与检验}
此处区分数据测量误差、模型近似误差和数值计算误差，并结合实际数据进行检验。

\subsection{敏感性分析}
此处说明参数扰动范围、评价指标及结论的稳定性。

% !TeX root = ../GMCM-main.tex
\section{模型的评价与推广}
\subsection{模型的优点}
此处结合计算结果，说明模型的有效性、解释性或计算效率。

\subsection{模型的不足}
此处说明模型依赖的条件、数据限制及尚未解决的问题。

\subsection{模型的推广}
此处说明可推广的场景、需要调整的条件与后续改进方向。

% !TeX root = ../GMCM-main.tex
\clearpage
% 保留 Paper 的 BibTeX 数据库及 gmcm 样式；请用真实引用替换示例。
文献引用排版示例：\cite{mittelbach_latex_2004,wright_latex3_2009,beeton_unicode_2008,vieth_experiences_2009}。
\bibliographystyle{gmcm}
\bibliography{reference}

% !TeX root = ../GMCM-main.tex
\clearpage
\begin{appendices}
  % 不重置页码，沿用 Paper 的附录环境。
  \section{代码程序}
  \renewcommand{\thesubsection}{\Alph{section}\thinskip.\thinskip\arabic{subsection}}
  下列程序改编自 CUMCM 的代码导入示例，仅用于展示外部源文件的排版。

  \subsection{Python 源程序}
  \lstinputlisting[style=Python,caption={三位数枚举示例}]{code/code_example.py}

  \subsection{MATLAB 源程序}
  \lstinputlisting[style=MATLAB,caption={等高线与梯度方向示例}]{code/code_example.m}

  \section{支撑材料}
  此处填写实际提交的程序、补充图表或其他支撑文件的名称及说明。
\end{appendices}

\end{document}
